\documentclass[11pt,a4paper]{article}
\usepackage[utf8]{inputenc}
\usepackage[T1]{fontenc}
\usepackage[english]{babel}
\usepackage[margin=2.5cm]{geometry}
\usepackage{booktabs}
\usepackage{longtable}
\usepackage{array}
\usepackage{xcolor}
\usepackage{tcolorbox}
\usepackage{enumitem}
\usepackage{hyperref}
\usepackage{amssymb}
\usepackage{fancyhdr}
\setlength{\headheight}{14pt}
\usepackage{titlesec}

% Colors
\definecolor{primaryblue}{HTML}{2c5aa0}
\definecolor{warningred}{HTML}{c62828}
\definecolor{successgreen}{HTML}{2a7a4b}
\definecolor{lightgray}{HTML}{f5f5f5}
\definecolor{darkgray}{HTML}{555555}

% Header/Footer
\pagestyle{fancy}
\fancyhf{}
\fancyhead[L]{\textcolor{darkgray}{\small Internal Document – School Administration}}
\fancyhead[R]{\textcolor{darkgray}{\small \thepage}}
\fancyfoot[C]{\textcolor{darkgray}{\small Confidential – For Administrative Use Only}}
\renewcommand{\headrulewidth}{0.4pt}
\renewcommand{\footrulewidth}{0.4pt}

% Section formatting
\titleformat{\section}{\large\bfseries\color{primaryblue}}{\thesection}{1em}{}
\titleformat{\subsection}{\normalsize\bfseries\color{primaryblue}}{\thesubsection}{1em}{}

% Box environments
\newtcolorbox{infobox}[1][]{
    colback=lightgray,
    colframe=primaryblue,
    title=#1,
    fonttitle=\bfseries,
    boxrule=1pt,
    arc=2pt
}

\newtcolorbox{warningbox}[1][]{
    colback=warningred!5,
    colframe=warningred,
    title=#1,
    fonttitle=\bfseries,
    boxrule=1.5pt,
    arc=2pt
}

\newtcolorbox{successbox}[1][]{
    colback=successgreen!5,
    colframe=successgreen,
    title=#1,
    fonttitle=\bfseries,
    boxrule=1pt,
    arc=2pt
}

\begin{document}

% Title
\begin{center}
    {\LARGE\bfseries\color{primaryblue} School Responsibilities \&\\[0.3em] Information Duties}\\[1em]
    {\large Regarding Regulatory Compliance (KLE, APVO M-V)}\\[0.5em]
    {\normalsize Internal Administrative Guide}\\[2em]
    {\small Document Version: January 2026}
\end{center}

\vspace{1em}

\begin{warningbox}[IMPORTANT NOTICE]
This document addresses the question of \textbf{who bears responsibility} when teachers are not properly informed about regulatory requirements, particularly regarding:
\begin{itemize}[nosep]
    \item KLE Sprechen (Complex Speaking Assessment)
    \item Teaching certification requirements (Lehrbefugnis)
    \item APVO M-V compliance in the qualification phase
\end{itemize}
\end{warningbox}

\vspace{1em}

%===================================================================
\section{Internal School Hierarchy of Responsibility}
%===================================================================

The following table outlines who is responsible for what within the school structure:

\begin{center}
\begin{tabular}{>{\bfseries}p{3.5cm}p{3.5cm}p{6.5cm}}
\toprule
\textbf{Level} & \textbf{Position} & \textbf{Responsibility} \\
\midrule
1. School Leadership & Principal (Schulleiter/in) &
    \begin{itemize}[nosep,leftmargin=*]
        \item Overall responsibility for APVO compliance
        \item Personnel matters
        \item Teaching certification verification
        \item Final accountability
    \end{itemize} \\[1em]
2. Upper Secondary & Coordinator (Oberstufenkoordinator) &
    \begin{itemize}[nosep,leftmargin=*]
        \item Operational Q-phase implementation
        \item KLE organization and scheduling
        \item Exam planning
        \item Informing department heads
    \end{itemize} \\[1em]
3. Department Level & Head of Languages (Fachkonferenzleitung) &
    \begin{itemize}[nosep,leftmargin=*]
        \item Subject-specific communication
        \item Passing requirements to teachers
        \item Coordinating within department
    \end{itemize} \\[1em]
4. Teacher Level & Subject Teacher (Fachlehrkraft) &
    \begin{itemize}[nosep,leftmargin=*]
        \item Implementation of assigned tasks
        \item \textbf{NOT} responsible for systemic issues
        \item \textbf{NOT} liable for missing certification
    \end{itemize} \\
\bottomrule
\end{tabular}
\end{center}

\subsection{Information Flow (Chain of Communication)}

\begin{verbatim}
    School Leadership (Principal)
           |
           +---> Upper Secondary Coordinator
           |              |
           |              +---> Head of Languages Department
           |                              |
           |                              +---> Spanish Teachers
           |
           +---> Direct: Personnel/certification matters
\end{verbatim}

\begin{infobox}[Key Principle]
Information about regulatory requirements must flow \textbf{downward} through official channels. Each level is responsible for ensuring the next level is properly informed.
\end{infobox}

%===================================================================
\section{Liability Analysis: Missed Information}
%===================================================================

\subsection{The Core Question}

\textbf{What happens when a teacher claims they were unaware of a requirement because they did not read an internal document?}

The answer depends on whether the school fulfilled its \textbf{duty to inform}.

\subsection{Burden of Proof}

\begin{warningbox}[The School Must Prove:]
\begin{enumerate}[nosep]
    \item The document was \textbf{delivered} to the teacher
    \item The teacher was \textbf{informed it was mandatory} reading
    \item \textbf{Reasonable time} was given to read and understand it
    \item The \textbf{importance} of the content was communicated
\end{enumerate}
\textbf{If the school cannot prove these points → The teacher is NOT liable.}
\end{warningbox}

%===================================================================
\section{Definition of Key Terms}
%===================================================================

\subsection{``Officially Distributed'' (Offiziell verteilt)}

\begin{successbox}[What Qualifies as Official Distribution]
A document is considered \textbf{officially distributed} when \textbf{ALL} of the following conditions are met:
\end{successbox}

\begin{center}
\begin{tabular}{>{\bfseries}p{4cm}p{9.5cm}}
\toprule
\textbf{Requirement} & \textbf{Description} \\
\midrule
1. Active Delivery &
    The document was actively sent/given to the teacher, not merely ``made available.'' \\[0.5em]

    \textcolor{successgreen}{\checkmark} & Email sent directly to teacher's official school address \\
    \textcolor{successgreen}{\checkmark} & Physical copy placed in teacher's mailbox \\
    \textcolor{successgreen}{\checkmark} & Handed out during a meeting \\[0.5em]
    \textcolor{warningred}{$\times$} & Uploaded to shared folder without notification \\
    \textcolor{warningred}{$\times$} & Posted on bulletin board \\
    \textcolor{warningred}{$\times$} & ``Available on the intranet'' \\
\midrule
2. Proof of Receipt &
    There is documentation that the teacher received it. \\[0.5em]

    \textcolor{successgreen}{\checkmark} & Signed distribution list \\
    \textcolor{successgreen}{\checkmark} & Email with read receipt confirmed \\
    \textcolor{successgreen}{\checkmark} & Minutes of meeting noting attendance \\[0.5em]
    \textcolor{warningred}{$\times$} & Email sent without confirmation \\
    \textcolor{warningred}{$\times$} & Assumed to be in mailbox \\
\midrule
3. Marked as Mandatory &
    The document was clearly labeled as required reading. \\[0.5em]

    \textcolor{successgreen}{\checkmark} & Subject line: ``MANDATORY: Please read by [date]'' \\
    \textcolor{successgreen}{\checkmark} & Cover page states ``Required reading for all Q-phase teachers'' \\
    \textcolor{successgreen}{\checkmark} & Verbal instruction in meeting (documented in minutes) \\[0.5em]
    \textcolor{warningred}{$\times$} & ``FYI'' or ``For your information'' \\
    \textcolor{warningred}{$\times$} & No indication of importance \\
\midrule
4. Reasonable Timeframe &
    Sufficient time was provided to read and act on the information. \\[0.5em]

    \textcolor{successgreen}{\checkmark} & Document provided weeks/months before deadline \\
    \textcolor{successgreen}{\checkmark} & Reminder sent before deadline \\[0.5em]
    \textcolor{warningred}{$\times$} & Document provided day before action required \\
    \textcolor{warningred}{$\times$} & No clear deadline communicated \\
\bottomrule
\end{tabular}
\end{center}

\subsection{``Bringschuld'' vs. ``Holschuld'' (Duty to Deliver vs. Duty to Seek)}

\begin{center}
\begin{tabular}{>{\bfseries}p{3cm}p{4.5cm}p{5.5cm}}
\toprule
\textbf{German Term} & \textbf{Meaning} & \textbf{Who Bears Risk} \\
\midrule
Bringschuld & Duty to deliver & \textbf{School} must actively push information to teachers \\[0.5em]
Holschuld & Duty to seek/fetch & \textbf{Teacher} must actively seek out information \\
\bottomrule
\end{tabular}
\end{center}

\begin{warningbox}[Critical Regulatory Matters = Bringschuld]
For matters involving:
\begin{itemize}[nosep]
    \item Legal requirements (APVO M-V, KLE)
    \item Examination regulations
    \item Certification requirements
    \item Student rights and obligations
\end{itemize}
The school has a \textbf{Bringschuld} (duty to deliver). The school leadership must \textbf{actively ensure} that all affected teachers are informed. It is \textbf{not sufficient} to say ``the information was available somewhere.''
\end{warningbox}

%===================================================================
\section{Liability Scenarios}
%===================================================================

\begin{center}
\begin{longtable}{p{6cm}p{3cm}p{4.5cm}}
\toprule
\textbf{Situation} & \textbf{Liability} & \textbf{Reasoning} \\
\midrule
\endhead

Document placed in shared folder without notification &
\textcolor{warningred}{\textbf{School's fault}} &
Teacher had no way of knowing it existed \\[0.5em]

Email sent, no read confirmation requested &
\textcolor{orange}{\textbf{Gray area}} &
School should have verified receipt \\[0.5em]

Email sent with read receipt confirmed &
\textcolor{successgreen}{\textbf{Teacher should have read}} &
Proof of receipt exists \\[0.5em]

Signed distribution list exists &
\textcolor{successgreen}{\textbf{Teacher liable}} &
Clear proof of receipt and acknowledgment \\[0.5em]

Mentioned in staff meeting (minutes exist, attendance documented) &
\textcolor{successgreen}{\textbf{Teacher should know}} &
Documented verbal communication \\[0.5em]

``It's on the intranet somewhere'' &
\textcolor{warningred}{\textbf{School's fault}} &
No active distribution, no notification \\[0.5em]

New teacher, no onboarding on regulations &
\textcolor{warningred}{\textbf{School's fault}} &
Duty to onboard new staff \\[0.5em]

Teacher absent during meeting, no follow-up &
\textcolor{warningred}{\textbf{School's fault}} &
Duty to inform absent staff separately \\[0.5em]

Annual reminder system in place, documented &
\textcolor{successgreen}{\textbf{Teacher should know}} &
Systematic approach to information \\

\bottomrule
\end{longtable}
\end{center}

%===================================================================
\section{The Teacher's Defense}
%===================================================================

\begin{infobox}[Valid Defense Statement]
A teacher can credibly defend themselves if they can truthfully state:

\begin{quote}
\textit{``I was never told this document existed. I was never informed that I needed to read it. I received no training or briefing on KLE requirements or the implications of my certification status for Q-phase teaching.''}
\end{quote}

If this statement is credible and the school cannot prove otherwise, \textbf{liability shifts entirely to school leadership}.
\end{infobox}

\subsection{What Teachers Should Document (Self-Protection)}

\begin{enumerate}
    \item Keep records of all official communications received
    \item Note dates of meetings attended and topics covered
    \item Request written confirmation of unclear verbal instructions
    \item Ask explicitly: ``Is there anything I need to read regarding [topic]?''
    \item Document when information was \textbf{not} provided
\end{enumerate}

%===================================================================
\section{Recommendations for School Administration}
%===================================================================

\subsection{Best Practices for Information Distribution}

\begin{enumerate}
    \item \textbf{Use a formal distribution system}
    \begin{itemize}
        \item Signed acknowledgment forms for critical documents
        \item Email with mandatory read receipts for digital distribution
        \item Documented handover during meetings
    \end{itemize}

    \item \textbf{Create an annual compliance calendar}
    \begin{itemize}
        \item Schedule regular reminders for recurring requirements (e.g., KLE deadlines)
        \item Document all reminders sent
    \end{itemize}

    \item \textbf{Implement onboarding procedures}
    \begin{itemize}
        \item Checklist for new teachers covering all regulatory requirements
        \item Signed acknowledgment of onboarding completion
    \end{itemize}

    \item \textbf{Follow up with absent staff}
    \begin{itemize}
        \item System to identify who missed critical meetings
        \item Documented follow-up communication
    \end{itemize}

    \item \textbf{Maintain a central compliance register}
    \begin{itemize}
        \item Record which teacher received which document and when
        \item Proof of distribution for audit purposes
    \end{itemize}
\end{enumerate}

%===================================================================
\section{Summary: Who Is Responsible?}
%===================================================================

\begin{center}
\begin{tabular}{p{6cm}p{7.5cm}}
\toprule
\textbf{Responsibility} & \textbf{Belongs To} \\
\midrule
Ensuring APVO M-V compliance & School Leadership \\
Organizing KLE assessments & Upper Secondary Coordinator \\
Informing teachers of requirements & School Leadership + Coordinator \\
Verifying teaching certifications & School Leadership (Personnel) \\
Reading officially distributed documents & Teacher (if properly distributed) \\
Seeking out undistributed information & \textbf{NOT the teacher's duty} \\
Consequences of systemic failures & School Leadership \\
\bottomrule
\end{tabular}
\end{center}

\vspace{1em}

\begin{warningbox}[Final Note]
\textbf{A teacher without proper teaching certification (Lehrbefugnis) who was assigned to teach a Q-phase course bears NO personal liability for:}
\begin{itemize}[nosep]
    \item The school's failure to conduct KLE assessments
    \item The mismatch between their qualifications and the assignment
    \item Regulatory non-compliance resulting from improper staffing
\end{itemize}
These are administrative decisions made \textbf{above the teacher's level} and remain the responsibility of school leadership.
\end{warningbox}

\vspace{2em}

\begin{center}
\rule{0.5\textwidth}{0.4pt}\\[1em]
{\small\textcolor{darkgray}{Document prepared for internal administrative use.\\
Not for distribution to students or external parties.}}
\end{center}

\end{document}
